\begin{abstract}
Large language model (LLM) inference engines are complex systems whose reliability affects correctness and service availability. This paper presents a code-centered empirical study of how bugs manifest, where faults reside, how developers repair them, and how repair scope evolves in six widely used open-source engines. We identify 50,121 pull requests related to bug fixing, submitted from the start of each project to March 2026. Among these, 12,707 were merged into their respective project codebases and manually verified as actual bug fixes, forming our study corpus. For each verified fix, we connect its observed symptom to the faulty code, repair operation, and patch scope. We find that 31.5\% of the verified fixes address bugs that directly affect correctness or availability. Symptom distributions also differ substantially across architectural layers. Maintenance is concentrated in both location and activity: the top 20\% of files account for 79.6\% of bug-fixing occurrences, while logic correction and adaptation make up 73.6\% of repair operations. From 2023 to early 2026, fixes increasingly span multiple files, while Hardware Adaptation occupies a growing share of bug-fixing activity. Together, these findings provide an empirical basis for improving how bugs are detected, localized, and repaired in LLM inference engines.
\end{abstract}
\begin{IEEEkeywords}
LLM inference engines, empirical study, software maintenance, software evolution.
\end{IEEEkeywords}
